\section{Introduction}\label{sec:intro}

A lattice L\'evy walk moves on $\Z$ in runs of alternating direction. The
run lengths are independent, and their tails are regularly varying with an
index $\alpha$ between $1$ and $2$, so a run has a finite mean and an
infinite variance. Let $A_N$ be the number of right steps among the first
$N$. The position of the walk after $N$ steps is $2A_N-N$, so the law
$P_N(k)=\Prob(A_N=k)$ on $\{0,\dots,N\}$ is the law of the walk read on its
natural lattice. L\'evy walks are a standard model of
superdiffusion~\cite{zdk2015}, and the same walk describes heat transport in
one-dimensional lattices~\cite{dsd2013,wlh2016} and the occupation time of an
alternating renewal process~\cite{godrecheluck2001,lamperti1958}. Its
profile has a central peak and two ballistic side
peaks~\cite{zdh2011,dzh2012}, and the physics literature describes the part
between them by an infinite density~\cite{rdhb2014,rdhb2014b}. We describe
$P_N(k)$ for every $k$, with proofs, with one exception that we state below.

\subsection{The picture}\label{sec:intro-picture}

Figure~\ref{fig:atlas} shows $P_N$ for $N=5000$ on a logarithmic scale.
Four pieces make up the row. We call the right steps plus slots.
\begin{itemize}[nosep]
\item The bulk. Near $pN$, where $p$ is the long-run fraction of plus
slots, the law is a stable density at the scale $B_N$ of the larger of the
two run-length fluctuations. Here $B_N$ grows like $N^{1/\alpha_*}$ up to a
slowly varying factor, with $\alpha_*$ the smaller index.
\item The big-jump zones. A count $k$ with $k-pN$ much larger than $B_N$
needs one plus run of macroscopic length, a giant run, and the law of that run
gives the bin its size, which is of order $T_\plus(N)$ on the ballistic range,
with $T_\plus$ the tail of a plus run. The giant run is either complete and lies
in the interior or the final run, censored at slot $N$, and the bin is the sum
of the two contributions.
\item The boundary layers. At $k=N-j$ with $j$ fixed, the giant run is the
final one, and the only question is where it started. It starts right after the $j$-th minus slot exactly when a
minus run ends there, so the layer is $T_\plus(N)$ times the renewal mass
function $u_\minus(j)$ of the minus runs.
\item The atoms. The bin $k=N$ is the event that the first run covers the
row, and its mass is the tail of the first run.
\end{itemize}
The first run is the preparation of the walk. It may be a fresh run, a run
seen from a typical time, or anything else, and it enters each of the four
pieces through its own criterion.

\begin{figure}[t]
\centering
\includegraphics[width=\textwidth]{atlas}
\caption{(a) The law $P_N(k)$ for $N=5000$, run lengths with
$f(n)\propto n^{-5/2}$ for both signs and a fresh plus start, on a
logarithmic scale, with the four pieces of Theorems~\ref{thm:bulk}
and~\ref{thm:bigjump}. The two ends of the horizontal axis are stretched.
(b) The position $j^*_N$ of the interior minimum near the plus end against
$N$, on logarithmic axes, with the law of Theorem~\ref{thm:antimode} and the
refined prediction with $C(j)$ from Theorem~\ref{thm:second-order}.
Section~\ref{sec:verify} describes the computations.}
\label{fig:atlas}
\end{figure}

\subsection{Results}\label{sec:intro-results}

We state the hypotheses in Section~\ref{sec:hypotheses}. They are regular
variation of the two tails with indices in $(1,2)$, written (RV), the local
form $f_\sgn(n)\sim\alpha_\sgn T_\sgn(n)/n$ of the mass functions, written
(LR), and aperiodicity, written (SPAN). The run laws may differ between the
two signs, so the two stable scales $b_{\plus,N}$ and $b_{\minus,N}$ may
differ, and their ratio may oscillate. Every theorem below holds for every
relation between the two scales. When the smaller scale is smaller than the
other by an unbounded factor, one window on the side of the sign with the
smaller scale is described only in part, as we say below. When the ratio of
the scales stays bounded, Theorem A covers that window. The model is
symmetric, so the side in question is the plus side when the plus scale is the
smaller one and the minus side otherwise. We write $s$ for the sum of the two mean run lengths and $p$
for the fraction of the plus mean.

\emph{The bulk.} Theorem~\ref{thm:bulk} says that for a fresh start, under
(RV) and (SPAN),
\[
 \sup_{k\in\Z}\bigl|B_NP_N(k)-h_N\bigl((k-pN)/B_N\bigr)\bigr|\to0,
 \qquad \sup_kP_N(k)\le\frac C{B_N},
\]
where $h_N$ is the density of $(1-p)\theta_{\plus,N}Z_\plus-p\,\theta_{\minus,N}Z_\minus$,
with $Z_\pm$ independent totally skewed stable variables. The supremum runs
over the whole row. On a window $|k-pN|\le KB_N$ the local limit also follows
from a strong renewal theorem of Berger, and what is new is the uniformity over
every bin under (RV) and (SPAN) alone. Outside the bulk the theorem says only that
every bin is
$o(1/B_N)$, and the next theorem says what the bins are.

\emph{Big jumps and the boundary layer.} Put $L_k=(k-pN)/(1-p)$, the length
of the plus run that carries a typical configuration to the count $k$.
Theorem~\ref{thm:bigjump} says that for a fresh plus start, under (RV) and
(SPAN) for both signs and (LR) for the plus sign,
\[
 P_N(k)=\Bigl[\gamma(N-k)\,f_\plus(\lfloor L_k\rfloor)+T_\plus(\lfloor L_k\rfloor)\,u_\minus(N-k)\Bigr](1+o(1))
\]
uniformly on $pN+K\lambda_NB_N\le k\le N$, in the sense that the error
tends to zero as $N\to\infty$ and then $K\to\infty$. Here $\gamma(j)$ is an
explicit finite sum with $\gamma(j)\sim js/\mu_\minus^2$, and
$\lambda_N=1+\log(B_N/b_{\plus,N})$, which equals $1$ when the plus scale is the
larger one. The first term is a complete giant
run and the second a censored final run. At $k=N-j$ with $j$ fixed the
theorem reduces to $P_N(N-j)/T_\plus(N)\to u_\minus(j)$, which needs only (RV)
for the plus sign. As $j$ grows, $u_\minus(j)\to1/\mu_\minus$ and the layer
joins the big-jump zone continuously. The junction with the bulk holds only
in the double limit. Remark~\ref{rem:no-single-formula} shows that neither
the bulk profile nor the big-jump kernel nor their sum approximates $P_N$
with relative error $o(1)$ over the plus side. When the plus scale is the
smaller one, on the window $KB_N\le k-pN\le K\lambda_NB_N$ we prove only the
bulk estimate and the lower bound $P_N(k)\ge(1-o(1))G_N(k)$, with $G_N$ the
bracket above, and we do not describe $P_N$ there. When the minus scale is the
smaller one the mirror window appears at the minus end. (LR) suffices for the
theorem, and Example~\ref{ex:rv-not-lr} shows that (RV) alone does not, since
the kernel reads $f_\plus$ at a single point. Whether some local condition
strictly weaker than (LR) suffices is open (Remark~\ref{rem:wr}). Read
at $k=vN$ the theorem gives the infinite density
\[
 \frac{P_N(\lfloor vN\rfloor)}{T_\plus(N)}\to
 \frac1{\mu_\minus}\Bigl[\alpha_\plus(1-\ell)\ell^{-\alpha_\plus-1}+\ell^{-\alpha_\plus}\Bigr],
 \qquad \ell=\frac{v-p}{1-p},
\]
with its censored-run term (Corollary~\ref{cor:infinite-density}).

\emph{Preparation.} Theorem~\ref{thm:preparation} describes every
preparation through the first-run decomposition of
Proposition~\ref{prop:first-run}. Write $\nu_\sgn(n)$ for the probability
that the first run has sign $\sgn$ and length $n$, and $Q_\sgn(n)$ for the
probability that it has sign $\sgn$ and length at least $n$. On every
macroscopic interior of the row the bulk approximation holds with no condition
on the preparation. It holds uniformly over $1\le k<N$ exactly when
$B_n(\nu_\plus(n)+\nu_\minus(n))\to0$, which we call the bulk criterion. The
preparation is invisible in the
boundary layer at the plus end exactly when $\nu_\plus(n)=o(T_\plus(n))$. On the
big-jump zone up to the end the same condition is sharp, while on its interior
only window averages of $\nu_\plus$ matter, and the window condition is sharp
under (LR) for both signs. At the plus end the conditions are nested,
since $Q_\plus(n)=o(T_\plus(n))$ implies $\nu_\plus(n)=o(T_\plus(n))$, which
implies $B_n\nu_\plus(n)\to0$, the plus half of the bulk criterion. When a
condition fails the theorem gives the exact correction, as a directional
defect in the bulk, and as an additive term in the big-jump zone and at the
end. The stationary preparation, which is
the equilibrated start of the physics literature, satisfies the bulk
criterion and fails the other two, and Corollary~\ref{cor:stationary} gives
its corrections in closed form. Corollary~\ref{cor:no-trim} shows that no
trimming of the row by $o(N)$ bins removes the dependence on the
preparation for every preparation at once.

\emph{The shape.} Theorem~\ref{thm:mode} compares the endpoint atoms with
the bulk for a preparation that satisfies the bulk criterion. An endpoint is
the global mode when its atom exceeds $H_N/B_N$ by a fixed factor, with $H_N$
the height of the profile, and the bulk wins when both atoms are smaller by a
fixed factor. For the stationary preparation
with equal indices this is the golden ratio threshold. The endpoints win for
$\alpha<\varphi=(1+\sqrt5)/2$, the bulk wins for $\alpha>\varphi$, and at
$\alpha=\varphi$ with pure power tails the bulk wins by the explicit factor
$1/\mathcal C(\varphi)=1.288\ldots$, with $\mathcal C$ as in~\eqref{eq:Calpha}.

Theorem~\ref{thm:antimode} locates the interior minimum between the bulk and
the plus end for a fresh plus start. Suppose that $u_\minus(j)>1/\mu_\minus$
for every $j\ge1$. This holds for every strictly log-convex law of the minus
runs. Then the minimum sits at $N-j^*_N$ with
\[
 T_\minus(j^*_N)\sim\frac{2\alpha_\plus s}N .
\]
So $j^*_N$ grows like the stable scale of the minus runs. At finite $N$ the
minimizer of $u_\minus(j)+C(j)/N$ predicts $j^*_N$ better, and $j^*_N$ is
asymptotic to it. Here $C(j)$ is the exact second-order coefficient of
Theorem~\ref{thm:second-order}. When $u_\minus$ dips below its limit the
minimum stays at bounded distance from the end. These results compare bins in
different zones. They do not say that the law is monotone between the
extrema, and the number of local extrema inside the layers is not determined
here.

\emph{Identification.} Proposition~\ref{thm:inverse} reads the run laws off the
ends. For a fresh plus start the atoms give the law of the plus runs directly,
since $P_N(N)=T_\plus(N)$. The law of the minus runs comes from the boundary
layer, since $P_N(N-j)/P_N(N)\to u_\minus(j)$ and the renewal equation inverts
$u_\minus$. The minus end cannot replace the atoms. Its ratios determine a
renewal mass function only up to a geometric thinning of the run law, and the
normalization of the law does not remove this freedom
(Remark~\ref{rem:thinning}).

\subsection{Related work}\label{sec:intro-related}

\emph{Infinite densities and big jumps.} Rebenshtok, Denisov, H\"anggi and
Barkai~\cite{rdhb2014,rdhb2014b} derived the infinite density of a L\'evy
walk from the Montroll and Weiss equation, formally and in continuous time.
For the two-state velocity model their density
\cite[Eqs.~(8) and~(10)]{rdhb2014} has the shape
$|v|^{-1-\alpha}[\alpha-(\alpha-1)|v|]$. A complete giant run alone would give
the linear term $-\alpha|v|$, and the censored final run adds $+|v|$, so their
formula already contains the censored run.
What is new in Corollary~\ref{cor:infinite-density} is a proof on the lattice
at fixed time, uniform up to the end, together with the boundary layer and the
preparation. In the physics model each flight chooses its direction
independently, and consecutive flights in the same direction merge into one
run. This gives the alternating walk with run laws thinned geometrically
with parameter $\frac12$, as in Remark~\ref{rem:thinning}. Vezzani, Barkai and
Burioni~\cite{vbb2019,vbb2020} derived the same tail from a single big jump.
Wang, Schulz, Deng and Barkai~\cite{wsdb2018} treat the occupation time of an
alternating renewal process with one law for both states and find a boundary
layer and an atom at the ends. They work in continuous time through a double
Laplace transform, with one law and a start of fair random sign, and they leave
the effect of the preparation on the rare fluctuations open. Here the layer is
proved on the lattice for two laws, uniformly up to $j=o(N)$, with a kernel
$\kappa$ that depends on the start (Theorem~\ref{thm:bigjump}(a) and~(c)).
Wang, H\"oll and Barkai~\cite{whb2020} treat the infinite-mean case.
Akimoto~\cite{akimoto2023} computes the moments of the occupation time for the
ordinary and the equilibrium start and shows that they differ for
$1<\alpha<2$. Theorem~\ref{thm:preparation} treats an arbitrary preparation bin by
bin and gives, for each zone, a sharp criterion for when the start is invisible
and the exact correction when it is not. The classical integral limit theorems for the sojourn time of an
alternating renewal process are due to Tak\'acs~\cite{takacs1957,takacs1959} and
Kesten~\cite{kesten1962}, and they are statements in distribution, while every
statement here concerns single bins. The
golden ratio appears in Wang, Li and H\"anggi~\cite{wlh2016} as the exponent
at which the side peaks and the central peak decay at the same rate, with the
edge exponent of Cipriani, Denisov and Politi~\cite{cdp2005}, and
Lee-Dadswell, Nickel and Gray~\cite{ldng2005} found the conduction exponent
$(3-\sqrt5)/2$ by a hydrodynamic argument. The threshold itself is the
exponent identity $\alpha-1=1/\alpha$. What Theorem~\ref{thm:mode} adds is the
comparison of single bins as a theorem and the constant $\mathcal C(\varphi)$.

\emph{Local limit theorems and renewal theory.} Doney~\cite{doney1991} proved
a bivariate local limit theorem in the stable domain, and
Berger~\cite{berger2019} proved multivariate local large deviations and strong
renewal theorems. His Theorem~3.3 applied to the walk of run lengths, together
with the fresh-plus identity of Proposition~\ref{prop:fresh-identity}, gives
the window $|k-pN|\le KB_N$ of Theorem~\ref{thm:bulk}, and Section~\ref{sec:bulk}
explains what the rest of the theorem adds and why it needs new work. Melbourne, P\`ene and Terhesiu~\cite{mpt2026} prove mixing local
limit theorems with multivariate stable limits on Young towers, which give one
bin at the stable scale with the first and the last run restricted to fixed
lengths. In the i.i.d.\ renewal setting Theorem~\ref{thm:bulk} removes these
restrictions. We use the local large deviation bound of Caravenna and
Doney~\cite{caravennadoney2019} (see also~\cite{melbourneterhesiu2022}) for sums
of a fixed number of variables. The local one-big-jump theorems for
sums of independent variables are due to Denisov, Dieker and
Shneer~\cite{dds2008} and Berger, Birkner and Yuan~\cite{bby2023}, and the one
under exactly the local condition (LR) is~\cite[Theorem~2.4]{berger2019ptrf}.
That theorem gives the probability that a fixed number of plus runs adds up to a
given large value. It does not give Theorem~\ref{thm:bigjump} directly, because
the number of runs before slot $N$ is random and tied to the run lengths, because
the big jump may be made by the final run, which the horizon cuts off and no
statement about complete sums sees, and because the kernel must hold uniformly up
to the endpoint, where only boundedly many minus slots remain. The proof of
Theorem~\ref{thm:bigjump} does not use these theorems. What is new there is
the occupation count at a fixed time, where a complete giant run enters with
the renewal density of its start and a censored run with the boundary kernel.
The second-order coefficient rests on the second-order renewal theorem of
Frenk~\cite{frenk1982}, see also Rogozin~\cite{rogozin1973} and
Doney~\cite{doney2023}, and the monotonicity of $u_\minus$ on de Bruijn and
Erd\H{o}s~\cite{debruijnerdos1953}.

\emph{The walk.} The walk is the directionally reinforced random walk of
Mauldin, Monticino and von Weizs\"acker~\cite{mmw1996}, with limit laws by
Horv\'ath and Shao~\cite{horvathshao1998}, and the persistent random walk of
C\'enac, Le Ny, de Loynes and Offret, who prove functional scaling limits,
a stable L\'evy process for $1<\alpha<2$~\cite{clnlo2018,clnlo2019}. Magdziarz,
Scheffler, Straka and \.Zebrowski~\cite{mssz2015} prove limit theorems for
L\'evy walks in continuous space and time. Theorem~\ref{thm:bulk} is the local
form of the marginal of these limits at one time, and the rest of the paper is
about the bins that the scaling limits do not see. Godr\`eche and
Luck~\cite{godrecheluck2001} give the transform of the occupation time, which is
the generating function of the fresh-plus identity, and Lamperti~\cite{lamperti1958}
proved the arcsine type law for infinite-mean runs.

\emph{This series.} For runs with geometric lengths the comparison of the
ends with the center is Problem~131 of Kagey's Open Problem
Collection~\cite{kagey131}, posed for a Galton board. The walk is then the
persistent random walk, whose crossings are treated in~\cite{paperA}, and
the elephant random walk with rare reversals, which has a Landau profile in
place of a stable one, is treated in~\cite{paperD}. The present paper is
the heavy-tailed member of the series.

\subsection{Method and organization}\label{sec:intro-organization}

Everything rests on the exact identities of Section~\ref{sec:setting}, which
write $P_N(k)$ through products of renewal kernels of the two signs and reduce
every preparation to the two fresh laws. The main results are five
theorems, lettered A to E. Sections~\ref{sec:bulk} to~\ref{sec:shape} prove
Theorems~\ref{thm:bulk} to~\ref{thm:antimode} in order,
Section~\ref{sec:inverse} proves Proposition~\ref{thm:inverse}, and
Section~\ref{sec:verify} describes the computations.
Appendix~\ref{app:classical} collects the classical inputs on regular variation,
the profile and sums of runs, and the bounds on the kernels of one sign.
